% ============================================================
% FUENTES: draft p.25--26.
% ============================================================
\section{Bifurcación transcrítica}

A diferencia de la bifurcación saddle-node, donde los puntos de equilibrio se crean o destruyen en pares, existen numerosos sistemas en la física, la química y la biología donde un punto de equilibrio específico persiste para todos los valores de los parámetros (por ejemplo, el estado de extinción $N=0$ en dinámica de poblaciones o el estado de reposo electromagnético en un láser). En tales casos, las ramas de equilibrio no desaparecen, sino que colisionan e intercambian su estabilidad. Este mecanismo se conoce como \textbf{bifurcación transcrítica}.

\subsection{Forma normal $\dot{x} = rx - x^{2}$ e intercambio de estabilidad}

La forma normal prototípica de la bifurcación transcrítica es:
\begin{equation}
    \dot{x} = rx - x^2 = x(r - x)
    \label{eq:transcritica_forma_normal}
\end{equation}
donde $r \in \mathbb{R}$ es el parámetro de control.

\paragraph{Puntos de equilibrio y estabilidad:}
Los puntos fijos se obtienen de $x(r - x) = 0$, rindiendo siempre dos soluciones:
\begin{equation}
    x_1^* = 0, \quad x_2^* = r
\end{equation}
Calculando la derivada del campo vectorial $f'(x) = r - 2x$:
\begin{itemize}
    \item En la rama trivial $x_1^* = 0$: $f'(0) = r$.
    \begin{itemize}
        \item Si $r < 0$: $f'(0) < 0 \implies x_1^* = 0$ es \textbf{Estable}.
        \item Si $r > 0$: $f'(0) > 0 \implies x_1^* = 0$ es \textbf{Inestable}.
    \end{itemize}
    \item En la rama móvil $x_2^* = r$: $f'(r) = r - 2r = -r$.
    \begin{itemize}
        \item Si $r < 0$: $f'(r) = -r > 0 \implies x_2^* = r$ es \textbf{Inestable}.
        \item Si $r > 0$: $f'(r) = -r < 0 \implies x_2^* = r$ es \textbf{Estable}.
    \end{itemize}
    \item En el valor crítico $r = 0$: Las dos ramas se cruzan en el origen $x^* = 0$, que se convierte en un nodo semi-estable $\dot{x} = -x^2$.
\end{itemize}

\begin{figure}[htbp]
\centering
\begin{tikzpicture}[>=Stealth, scale=0.82]
    % Panel 1: r < 0
    \begin{scope}[shift={(-6.2, 3.2)}]
        \draw[->, thick, black!60] (-2.2,0) -- (2.2,0) node[right, font=\tiny] {$x$};
        \draw[->, thick, black!60] (0,-1.5) -- (0,1.8) node[above, font=\tiny] {$\dot{x}$};
        \draw[thick, black!85, domain=-1.6:0.8] plot (\x, {-1.0*\x - \x*\x});
        \draw[black, fill=white, thick] (-1.0, 0) circle (2pt) node[below=2pt, font=\tiny] {$x^*=r$};
        \filldraw[black] (0, 0) circle (2pt) node[above right=1pt, font=\tiny] {$x^*=0$};
        \draw[->, thick, black!90] (-1.6,0) -- (-1.2,0);
        \draw[->, thick, black!90] (-0.3,0) -- (-0.7,0);
        \draw[->, thick, black!90] (0.8,0) -- (0.3,0);
        \node[font=\footnotesize, align=center] at (0,-1.8) {\textbf{(a)} $r < 0$ ($x=0$ estable)};
    \end{scope}

    % Panel 2: r = 0
    \begin{scope}[shift={(0, 3.2)}]
        \draw[->, thick, black!60] (-2.2,0) -- (2.2,0) node[right, font=\tiny] {$x$};
        \draw[->, thick, black!60] (0,-1.5) -- (0,1.8) node[above, font=\tiny] {$\dot{x}$};
        \draw[thick, black!85, domain=-1.2:1.2] plot (\x, {-\x*\x});
        \filldraw[black!60] (0, 0) circle (2pt) node[above=2pt, font=\tiny] {$x^*=0$};
        \draw[->, thick, black!90] (-1.2,0) -- (-0.4,0);
        \draw[->, thick, black!90] (1.2,0) -- (0.4,0);
        \node[font=\footnotesize, align=center] at (0,-1.8) {\textbf{(b)} $r = 0$ (semi-estable)};
    \end{scope}

    % Panel 3: r > 0
    \begin{scope}[shift={(6.2, 3.2)}]
        \draw[->, thick, black!60] (-2.2,0) -- (2.2,0) node[right, font=\tiny] {$x$};
        \draw[->, thick, black!60] (0,-1.5) -- (0,1.8) node[above, font=\tiny] {$\dot{x}$};
        \draw[thick, black!85, domain=-0.8:1.6] plot (\x, {1.0*\x - \x*\x});
        \draw[black, fill=white, thick] (0, 0) circle (2pt) node[below left=1pt, font=\tiny] {$x^*=0$};
        \filldraw[black] (1.0, 0) circle (2pt) node[above=2pt, font=\tiny] {$x^*=r$};
        \draw[->, thick, black!90] (-0.8,0) -- (-0.3,0);
        \draw[->, thick, black!90] (0.3,0) -- (0.7,0);
        \draw[->, thick, black!90] (1.6,0) -- (1.2,0);
        \node[font=\footnotesize, align=center] at (0,-1.8) {\textbf{(c)} $r > 0$ ($x=r$ estable)};
    \end{scope}

    % Panel 4: Diagrama de bifurcación (r vs x*)
    \begin{scope}[shift={(0, -2.8)}]
        \draw[->, thick, black!60] (-3.0,0) -- (3.0,0) node[right, font=\small] {$r$};
        \draw[->, thick, black!60] (0,-2.2) -- (0,2.2) node[above, font=\small] {$x^*$};
        
        % Rama trivial x* = 0
        \draw[very thick, black!85] (-2.5,0) -- (0,0); % estable para r < 0
        \draw[very thick, dashed, black!75] (0,0) -- (2.5,0); % inestable para r > 0
        
        % Rama móvil x* = r
        \draw[very thick, dashed, black!75] (-2.0,-2.0) -- (0,0); % inestable para r < 0
        \draw[very thick, black!85] (0,0) -- (2.0,2.0); % estable para r > 0
        
        \filldraw[black] (0,0) circle (2.5pt) node[below right=2pt, font=\footnotesize] {Bifurcación $(0,0)$};
        \node[above, font=\footnotesize] at (-1.5, 0.1) {Estable ($x^*=0$)};
        \node[below, font=\footnotesize] at (1.5, -0.1) {Inestable ($x^*=0$)};
        \node[above, font=\footnotesize] at (1.2, 1.3) {Estable ($x^*=r$)};
        \node[font=\bfseries\small] at (0, -2.6) {\textbf{(d)} Intercambio de estabilidad (\emph{exchange of stabilities})};
    \end{scope}
\end{tikzpicture}
\caption{Bifurcación transcrítica en $\dot{x} = rx - x^2$. Las ramas de equilibrio $x^*=0$ y $x^*=r$ coexisten para todo $r$ e intercambian su estabilidad al colisionar en el punto crítico $r=0$.}
\label{fig:bifurcacion_transcritica}
\end{figure}

\subsection{Ejemplo no polinomial y condición $ab = 1$}

Consideremos el siguiente flujo unidimensional no polinomial:
\begin{equation}
    \dot{x} = x(1 - x^2) - a(1 - e^{-bx})
    \label{eq:ejemplo_transcritica_nopol}
\end{equation}
con parámetros positivos $a > 0$ y $b > 0$.

Notemos en primer lugar que $x^* = 0$ es un punto de equilibrio para cualquier par $(a, b)$, ya que $0(1 - 0) - a(1 - e^0) = 0$.

\paragraph{Aproximación en serie de Taylor.}
Para analizar la dinámica cerca del origen ($|x| \ll 1$), expandimos la exponencial $e^{-bx}$ en serie de Maclaurin:
\begin{equation}
    e^{-bx} = 1 - bx + \frac{1}{2}b^2 x^2 - \frac{1}{6}b^3 x^3 + \mathcal{O}(x^4)
\end{equation}
Sustituyendo en la ecuación diferencial \eqref{eq:ejemplo_transcritica_nopol}:
\begin{align}
    \dot{x} &= x - x^3 - a\left[1 - \left(1 - bx + \frac{1}{2}b^2 x^2 - \frac{1}{6}b^3 x^3 + \mathcal{O}(x^4)\right)\right] \nonumber \\
    &= x - x^3 - a\left(bx - \frac{1}{2}b^2 x^2 + \mathcal{O}(x^3)\right) \nonumber \\
    &= (1 - ab)x + \frac{1}{2}ab^2 x^2 + \mathcal{O}(x^3)
    \label{eq:expansion_transcritica_ab}
\end{align}

\begin{figure}[htbp]
\centering
\begin{tikzpicture}[>=Stealth, scale=0.9]
    \draw[->, thick, black!60] (0,0) -- (4.5,0) node[right, font=\small] {$a$};
    \draw[->, thick, black!60] (0,0) -- (0,3.8) node[above, font=\small] {$b$};
    
    % Hipérbola ab = 1 => b = 1/a
    \draw[very thick, black!85, domain=0.3:4.0, samples=60] plot (\x, {1.0/\x}) node[above right, font=\footnotesize] {$ab = 1$ (curva de bifurcación)};
    
    \node[font=\footnotesize, black!75, align=center] at (2.8, 2.2) {Régimen $ab > 1$:\\Origen inestable ($f'(0) < 0$),\\nace rama no trivial estable $x^* > 0$};
    \node[font=\footnotesize, black!75, align=center] at (1.0, 0.4) {Régimen $ab < 1$:\\Origen $x^*=0$ estable};
    
    \filldraw[black] (1.0, 1.0) circle (2pt) node[below left, font=\tiny] {$(1,1)$};
\end{tikzpicture}
\caption{Curva de bifurcación transcrítica en el espacio de parámetros $(a, b)$. La hipérbola $ab = 1$ separa las regiones de estabilidad del estado trivial.}
\label{fig:curva_bifurcacion_ab}
\end{figure}

\paragraph{Condición de bifurcación y rama no nula.}
Comparando con la forma normal $\dot{x} = R x + c x^2$, la bifurcación transcrítica ocurre cuando el coeficiente lineal se anula:
\begin{equation}
    1 - ab = 0 \iff \boxed{ab = 1}
\end{equation}
Esta relación define una hipérbola crítica en el primer cuadrante del plano de parámetros $(a, b)$ (véase la \cref{fig:curva_bifurcacion_ab}).

Para encontrar la fórmula aproximada del punto fijo que bifurca desde $x=0$ en las inmediaciones de la curva ($ab \approx 1$):
\begin{equation}
    (1 - ab)x^* + \frac{ab^2}{2}(x^*)^2 \approx 0 \implies x^* \left[(1 - ab) + \frac{ab^2}{2}x^*\right] \approx 0
\end{equation}
Excluyendo la solución trivial $x^* = 0$, obtenemos la rama bifurcada:
\begin{equation}
    x^* \approx \frac{2(ab - 1)}{ab^2}
    \label{eq:rama_bifurcada_ab}
\end{equation}
Esta expresión es válida localmente para $|x^*| \ll 1$ y $|ab - 1| \ll 1$.

